\documentclass[12pt,a4paper]{article}
\usepackage[utf8]{inputenc}
\usepackage{amsmath,amssymb,amsthm}
\usepackage{geometry}
\geometry{margin=1in}

\newtheorem{theorem}{Theorem}
\newtheorem{lemma}[theorem]{Lemma}
\newtheorem{corollary}[theorem]{Corollary}

\title{Problem 69: Totient Maximum}
\author{Project Euler Solution}
\date{}

\begin{document}
\maketitle

\section{Problem Statement}

Find the value of $n \le 1\,000\,000$ for which $n / \varphi(n)$ is a maximum, where $\varphi$ is Euler's totient function.

\section{Mathematical Analysis}

\begin{theorem}[Ratio Formula]
For $n = p_1^{a_1} p_2^{a_2} \cdots p_k^{a_k}$, the ratio $n/\varphi(n)$ depends only on the set of distinct prime divisors:
\[
\frac{n}{\varphi(n)} = \prod_{p \mid n} \frac{p}{p - 1}.
\]
\end{theorem}

\begin{proof}
By Euler's product formula, $\varphi(n) = n \prod_{p \mid n}(1 - 1/p)$. Dividing:
\[
\frac{n}{\varphi(n)} = \prod_{p \mid n} \frac{1}{1 - 1/p} = \prod_{p \mid n} \frac{p}{p-1}.
\]
The exponents $a_i$ cancel, so the ratio depends only on $\{p_1, \ldots, p_k\}$.
\end{proof}

\begin{theorem}[Monotonicity]
\begin{enumerate}
    \item Adding any prime $q$ to the set of divisors multiplies $n/\varphi(n)$ by $q/(q-1) > 1$.
    \item The factor $p/(p-1) = 1 + 1/(p-1)$ is strictly decreasing in $p$.
\end{enumerate}
\end{theorem}

\begin{proof}
Statement (1) is immediate since $q/(q-1) > 1$ for all primes $q \geq 2$. Statement (2) follows because $1/(p-1)$ is strictly decreasing.
\end{proof}

\begin{theorem}[Primorial Optimality]
Among all $n \leq N$, the ratio $n/\varphi(n)$ is maximized when $n$ is the largest primorial not exceeding~$N$:
\[
n^* = \prod_{\substack{i=1}}^{k} p_i, \quad \text{where } p_1 p_2 \cdots p_k \leq N < p_1 p_2 \cdots p_{k+1}.
\]
\end{theorem}

\begin{proof}
We prove optimality in three steps:

\textbf{Step 1 (Squarefree).} If $n$ is not squarefree, say $n = p^2 m$ with $p \nmid m$, then $n' = pm$ satisfies $n' < n$ and $n'/\varphi(n') = n/\varphi(n)$ by Theorem~1. The freed factor of~$p$ can be used to include a new prime, strictly increasing the ratio by Theorem~2(1). Hence the optimizer is squarefree.

\textbf{Step 2 (Consecutive primes).} If $n = p_{j_1} \cdots p_{j_k}$ with some $j_i > i$, replacing $p_{j_i}$ by $p_i$ (a smaller prime) does not increase $n$ and strictly increases the ratio by Theorem~2(2). Hence the optimizer uses consecutive smallest primes.

\textbf{Step 3 (Maximal count).} By Theorem~2(1), each additional prime strictly increases the ratio, so we include as many consecutive primes as the bound allows.
\end{proof}

\subsection{Computation}

\[
2 \times 3 \times 5 \times 7 \times 11 \times 13 \times 17 = 510510 \le 1\,000\,000
\]
\[
2 \times 3 \times 5 \times 7 \times 11 \times 13 \times 17 \times 19 = 9\,699\,690 > 1\,000\,000
\]

\begin{corollary}
The answer is $n = 510510$ with:
\[
\frac{510510}{\varphi(510510)} = \frac{2}{1} \cdot \frac{3}{2} \cdot \frac{5}{4} \cdot \frac{7}{6} \cdot \frac{11}{10} \cdot \frac{13}{12} \cdot \frac{17}{16} \approx 5.214
\]
\end{corollary}

\section{Complexity}

\begin{itemize}
    \item \textbf{Time}: $O(\log \log N)$ multiplications --- we multiply consecutive primes until exceeding~$N$.
    \item \textbf{Space}: $O(1)$.
\end{itemize}

\section{Result}

The answer is $\boxed{510510}$.

\section{Answer}

\[
\boxed{510510}
\]

\end{document}
